\scanpage{364}\mspage{280}
Si on pouvait trouver une sous-catégorie
\inserted{(équiv. à une petite catégorie)}
$\unsure{\mathscr H|\mathcal H}_\pi$ de $\mathscr H=\underline{\mathrm{Hom}}_{\mathrm{acc}}(M,N)$, \inserted{\unsure{pour} un cardinal $\pi\geq\pi_0$,} telle que tt
objet $F$ de $\mathscr H$ s'écrive $\underset{\unsure{\mathbb I|I|J}}{\unsure{\underline{\lim}|\underrightarrow{\lim}}}\,F_i$, les
$F_i$ dans $\mathscr H_\pi$, $\mathbb I$ grd devant $\pi$, alors
\unsure{\{|\deleted{à}} si les objets de $\mathscr H_\pi$ sont \deleted{\unsure{$\pi$-access}}\inserted{des}
foncteurs $\pi'$-accessibles de $M$ dans $N$\unsure{\}}, 
il en serait de $=$ de \unsure{tout|tous} $F$ dans $\mathscr H$.
Donc il existerait un cardinal \inserted{infini} $\pi'$
tel que tout foncteur accessible de
$M$ dans $N$ soit $\pi'$-accessible. Mais
sauf erreur, si \inserted{p.ex.} $M=N=\mathrm{Ens}$, il n'existe
pas de tel cardinal. Je ne veux
pas m'arrêter ici sur cette question,
et faire un contre-exemple.
% archive sheet number: 363

